% !TEX root = main.tex
\chapter{Fraïssé Classes and Fraïssé Limits}

This chapter develops the Fraïssé-theoretic background needed for the rest of the thesis. We introduce the basic notions of age, hereditary property, joint embedding property, and amalgamation property, and we explain how these conditions characterize classes of finite structures arising from countable ultrahomogeneous structures. We then state and prove Fraïssé's theorem, which yields the existence and uniqueness of the Fraïssé limit of a Fraïssé class. These constructions provide the model-theoretic foundation for the later study of automorphism groups, Ramsey properties, and universal minimal flows.

\section{Preliminaries on structures and embeddings}

Throughout this chapter, \(L\) denotes a signature. An \(L\)-structure \(A\) consists of a set \(\mathrm{dom}(A)\) together with interpretations of the constant, relation, and function symbols of \(L\). If \(c\) is a constant symbol, \(R\) an \(n\)-ary relation symbol, and \(F\) an \(n\)-ary function symbol, their interpretations in \(A\) are denoted by
\[
    c^A\in \mathrm{dom}(A),\qquad R^A\subseteq \mathrm{dom}(A)^n,\qquad F^A:\mathrm{dom}(A)^n\to \mathrm{dom}(A).
\]
We write \(|A|\) for \(|\mathrm{dom}(A)|\). Constants may, when convenient, be regarded as \(0\)-ary function symbols.

If \(A\) and \(B\) are \(L\)-structures, a map \(f:\mathrm{dom}(A)\to \mathrm{dom}(B)\) is a \emph{homomorphism} if
\begin{enumerate}
    \item \(f(c^A)=c^B\) for every constant symbol \(c\),
    \item whenever \(\bar a\in R^A\), one has \(f(\bar a)\in R^B\),
    \item \(f(F^A(\bar a))=F^B(f(\bar a))\) for every tuple \(\bar a\).
\end{enumerate}
An \emph{embedding} is an injective homomorphism which also reflects relations, that is,
\[
    \bar a\in R^A \iff f(\bar a)\in R^B
\]
for every relation symbol \(R\). A bijective embedding is an \emph{isomorphism}. We write \(A\cong B\) if \(A\) and \(B\) are isomorphic. An \emph{automorphism} of \(A\) is an isomorphism from \(A\) onto itself.


\begin{lemma}\label{lem:isomorphism-basics}
    Let \(A,B,C\) be \(L\)-structures.
    \begin{enumerate}
        \item[(i)] The identity map \(1_A\colon A\to A\) is an isomorphism.
        \item[(ii)] If \(f\colon A\to B\) is an isomorphism, then \(f^{-1}\colon B\to A\) is an isomorphism.
        \item[(iii)] If \(f\colon A\to B\) and \(g\colon B\to C\) are isomorphisms, then \(g\circ f\colon A\to C\) is an isomorphism.
    \end{enumerate}
\end{lemma}

\begin{proof}
    Part (i) is immediate. For (ii), since \(f\) is a bijective embedding, its inverse is again a bijection preserving constants, relations, and functions, and hence is an isomorphism. For (iii), by composition of embeddings, \(g\circ f\) is an embedding, and since it is also a composition of bijections, it is bijective. Therefore \(g\circ f\) is an isomorphism.
\end{proof}

\begin{definition}[Substructure]
    Let \(A\) and \(B\) be \(L\)-structures with \(\mathrm{dom}(A)\subseteq \mathrm{dom}(B)\). We say that \(A\) is a \emph{substructure} of \(B\), and write \(A\subseteq B\), if the inclusion map is an embedding. Equivalently,
    \[
        c^A=c^B,\qquad R^A=R^B\cap \mathrm{dom}(A)^n,\qquad
        F^A=F^B\!\restriction_{\mathrm{dom}(A)^n}
    \]
    for all constant, relation, and function symbols of \(L\).
\end{definition}

\begin{lemma}\label{lem:substructure-characterization}
    Let \(B\) be an \(L\)-structure and let \(X\subseteq \mathrm{dom}(B)\). The following are equivalent:
    \begin{enumerate}
        \item[(i)] \(X=\mathrm{dom}(A)\) for some substructure \(A\subseteq B\);
        \item[(ii)] \(X\) contains all constants of \(B\) and is closed under all functions of \(B\).
    \end{enumerate}
    If these conditions hold, the substructure \(A\) is uniquely determined.
\end{lemma}

\begin{proof}
    If \(A\subseteq B\), then \(\mathrm{dom}(A)\) plainly contains the constants and is closed under the functions of \(B\). Conversely, if \(X\subseteq \mathrm{dom}(B)\) satisfies (ii), define an \(L\)-structure \(A\) on \(X\) by
    \[
        c^A=c^B,\qquad R^A=R^B\cap X^n,\qquad F^A=F^B\!\restriction_{X^n}.
    \]
    Then the inclusion \(A\hookrightarrow B\) is an embedding, so \(A\subseteq B\). Uniqueness is immediate from the definition of substructure.
\end{proof}

\begin{definition}[Generated substructure]
    Let \(B\) be an \(L\)-structure and \(Y\subseteq \mathrm{dom}(B)\). By Lemma~\ref{lem:substructure-characterization}, there is a unique smallest substructure of \(B\) whose domain contains \(Y\). It is denoted by
    \[
        \langle Y\rangle_B
    \]
    and is called the \emph{substructure generated by \(Y\)}. When the ambient structure is clear, we simply write \(\langle Y\rangle\). A structure is \emph{finitely generated} if it is generated by a finite subset of its domain.
\end{definition}

The following observation is used repeatedly: any embedding may be viewed as an inclusion after replacing the codomain by an isomorphic copy.

\begin{lemma}\label{lem:embedding-as-inclusion}
    Let \(f:A\hookrightarrow C\) be an embedding of \(L\)-structures. Then there exist an \(L\)-structure \(B\) and an isomorphism \(h:B\to C\) such that \(A\subseteq B\) and
    \[
        f=h\!\restriction_A .
    \]
\end{lemma}

\begin{proof}
    Let
    \[
        \mathrm{dom}(B):=\mathrm{dom}(A)\,\dot\cup\,\bigl(\mathrm{dom}(C)\setminus f(\mathrm{dom}(A))\bigr),
    \]
    and define a bijection \(\sigma:\mathrm{dom}(B)\to \mathrm{dom}(C)\) by
    \[
        \sigma(a)=f(a)\quad(a\in A),\qquad \sigma(x)=x\quad\bigl(x\in \mathrm{dom}(C)\setminus f(\mathrm{dom}(A))\bigr).
    \]
    Transport the \(L\)-structure of \(C\) to \(\mathrm{dom}(B)\) via \(\sigma\). Then \(\sigma\) becomes an isomorphism \(h:B\to C\). Since \(f\) is an embedding, the induced structure on \(\mathrm{dom}(A)\) is exactly the original structure \(A\), so \(A\subseteq B\) and \(f=h\!\restriction_A\).
\end{proof}

For the examples relevant to Fraïssé theory, the language is usually finite and relational. In that case finitely generated structures are automatically finite.

\begin{proposition}\label{prop:finite-relational-language-finiteness}
    Let \(L\) be a finite relational signature.
    \begin{enumerate}
        \item Every finitely generated \(L\)-structure is finite.
        \item For each \(n\geq 1\), there are only finitely many \(L\)-structures of cardinality \(n\), up to isomorphism.
    \end{enumerate}
\end{proposition}

\begin{proof}
    Since \(L\) has no function symbols and no constants, the substructure generated by a set is just that set itself. Hence every finitely generated \(L\)-structure is finite.

    Now fix \(n\geq 1\). Every \(L\)-structure of size \(n\) is isomorphic to one with domain \(\{1,\dots,n\}\). If \(R_1,\dots,R_k\) are the relation symbols of \(L\), of arities \(m_1,\dots,m_k\), then each \(R_i\) may be interpreted arbitrarily as a subset of \(\{1,\dots,n\}^{m_i}\). There are only finitely many such choices, so only finitely many \(L\)-structures on \(\{1,\dots,n\}\), and therefore only finitely many isomorphism types of size \(n\).
\end{proof}


\section{Ages, HP, and JEP}

We now introduce the first invariant associated with a structure.

\begin{definition}[Age]
    Let \(L\) be a signature and let \(D\) be an \(L\)-structure. The \emph{age} of \(D\), denoted \(\Age(D)\), is the class of all finitely generated \(L\)-structures that embed into \(D\). Equivalently, \(\Age(D)\) is the class of all finitely generated substructures of \(D\), taken up to isomorphism.
\end{definition}

In this thesis we regard ages up to isomorphism. Thus to say that \(D\) has countable age means that there are only countably many isomorphism types of finitely generated substructures of \(D\).

The following two properties are immediate for every age.

\begin{definition}
    Let \(\mathcal K\) be a class of finitely generated \(L\)-structures.
    \begin{enumerate}
        \item[\textup{(HP)}] \(\mathcal K\) has the \emph{hereditary property} if whenever \(A\in\mathcal K\) and \(B\) is a finitely generated substructure of \(A\), then \(B\in\mathcal K\) up to isomorphism.
        \item[\textup{(JEP)}] \(\mathcal K\) has the \emph{joint embedding property} if for all \(A,B\in\mathcal K\), there exists \(C\in\mathcal K\) such that both \(A\) and \(B\) embed into \(C\).
    \end{enumerate}
\end{definition}

\begin{proposition}\label{prop:age-hp-jep}
    Let \(D\) be an \(L\)-structure. Then \(\Age(D)\) is nonempty and satisfies \emph{(HP)} and \emph{(JEP)}.
\end{proposition}

\begin{proof}
    Choose any \(a\in D\). Then the substructure \(\langle a\rangle\) generated by \(a\) is finitely generated and embeds into \(D\), so \(\Age(D)\neq\emptyset\).

    To prove \emph{(HP)}, let \(A\in\Age(D)\), and let \(B\subseteq A\) be finitely generated. Since \(A\) embeds into \(D\), the restriction of any embedding \(A\hookrightarrow D\) gives an embedding \(B\hookrightarrow D\). Hence \(B\in\Age(D)\).

    To prove \emph{(JEP)}, let \(A,B\in\Age(D)\). Choose embeddings \(e_A:A\hookrightarrow D\) and \(e_B:B\hookrightarrow D\). Let
    \[
        C:=\langle e_A(A)\cup e_B(B)\rangle \subseteq D.
    \]
    Since \(A\) and \(B\) are finitely generated, the structure \(C\) is finitely generated. Thus \(C\in\Age(D)\), and both \(A\) and \(B\) embed into \(C\).
\end{proof}

Fraïssé proved a converse: under mild countability assumptions, these two properties already characterize the ages of countable structures.

\begin{proposition}\label{prop:hp-jep-age}
    Let \(L\) be a signature, and let \(\mathcal K\) be a nonempty class of finitely generated \(L\)-structures, countable up to isomorphism. Suppose that \(\mathcal K\) satisfies \emph{(HP)} and \emph{(JEP)}. Then there exists a countable \(L\)-structure \(C\) such that
    \[
        \mathcal K=\Age(C).
    \]
\end{proposition}

\begin{proof}
    Choose a sequence \((A_i)_{i<\omega}\) listing all isomorphism types in \(\mathcal K\), with each type appearing at least once. We construct a chain
    \[
        B_0\subseteq B_1\subseteq B_2\subseteq \cdots
    \]
    such that each \(B_i\in\mathcal K\) up to isomorphism and \(A_i\) embeds into \(B_{i+1}\).

    Set \(B_0=A_0\). Suppose \(B_i\) has been constructed. By \emph{(JEP)}, there is some \(C_i\in\mathcal K\) into which both \(B_i\) and \(A_{i+1}\) embed. By Lemma~\ref{lem:embedding-as-inclusion}, after replacing \(C_i\) by an isomorphic copy, we may assume that \(B_i\subseteq C_i\). Set \(B_{i+1}:=C_i\). This completes the construction.

    Now let
    \[
        C:=\bigcup_{i<\omega} B_i .
    \]
    Since \(C\) is a union of countably many finitely generated structures, it is countable.

    By construction, every member of \(\mathcal K\) embeds into \(C\), so \(\mathcal K\subseteq \Age(C)\). Conversely, let \(A\) be a finitely generated substructure of \(C\). A finite generating set for \(A\) is contained in some \(B_i\), hence \(A\subseteq B_i\). Since \(B_i\in\mathcal K\) up to isomorphism and \(\mathcal K\) satisfies \emph{(HP)}, it follows that \(A\in\mathcal K\) up to isomorphism. Therefore \(\Age(C)\subseteq \mathcal K\).

    Hence \(\Age(C)=\mathcal K\).
\end{proof}

\begin{remark}
    If \(L\) is a finite relational signature, then every finitely generated \(L\)-structure is finite. In that case the countability assumption on \(\mathcal K\) is often easy to verify by Proposition~\ref{prop:finite-relational-language-finiteness}.
\end{remark}


\section{A motivating example: finite linear orders}

The preceding discussion shows that every age satisfies \emph{(HP)} and \emph{(JEP)}. A natural question is whether these conditions are strong enough to determine a distinguished countable structure. The example of linear orders shows that they are not.

\begin{proposition}\label{prop:age-infinite-linear-order}
    Let \(L=\{<\}\), and let \((M,<)\) be an infinite linear order. Then
    \[
        \Age(M)=\{\text{finite linear orders}\}.
    \]
    In particular, any two infinite linear orders have the same age.
\end{proposition}

\begin{proof}
    Since the language \(L=\{<\}\) is purely relational, every finitely generated substructure of \(M\) is simply the induced substructure on a finite subset of \(M\). The restriction of \(<\) to a finite subset is again a finite linear order. Thus every member of \(\Age(M)\) is a finite linear order.

    Conversely, let \((F,<)\) be a finite linear order, say with
    \[
        F=\{a_1,\dots,a_n\}, \qquad a_1<\cdots<a_n .
    \]
    Since \(M\) is infinite, we may choose distinct elements \(b_1,\dots,b_n\in M\), and after reordering assume that
    \[
        b_1<\cdots<b_n .
    \]
    Then the map \(e\colon F\to M\) given by \(e(a_i)=b_i\) is an embedding. Hence every finite linear order belongs to \(\Age(M)\), and therefore
    \[
        \Age(M)=\{\text{finite linear orders}\}.
    \]
\end{proof}

Let \(\mathcal K_{<}\) denote the class of all finite linear orders, up to isomorphism. By Proposition~\ref{prop:age-infinite-linear-order},
\[
    \Age(\mathbb Q,<)=\Age(\mathbb Z,<)=\mathcal K_{<}.
\]
Moreover, \(\mathcal K_{<}\) satisfies \emph{(HP)} and \emph{(JEP)}: \emph{(HP)} is immediate, and \emph{(JEP)} holds because any two finite linear orders embed into a sufficiently long finite linear order.

Thus age, together with \emph{(HP)} and \emph{(JEP)}, does not distinguish the countable dense linear order without endpoints from the order of the integers. In particular, these conditions do not yet explain why the finite linear orders ``tend'' to \((\mathbb Q,<)\) rather than to \((\mathbb Z,<)\).

This shows that one further condition is needed. The next section introduces the \emph{amalgamation property}, which is precisely the additional condition singled out by Fraïssé.

\section{The amalgamation property and Fraïssé classes}

The previous section shows that age, together with \emph{(HP)} and \emph{(JEP)}, does not suffice to distinguish \((\mathbb Q,<)\) from \((\mathbb Z,<)\). Fraïssé's additional requirement is the amalgamation property.

\begin{definition}[Amalgamation property]
    Let \(\mathcal K\) be a class of finitely generated \(L\)-structures. We say that \(\mathcal K\) has the \emph{amalgamation property}, or \emph{(AP)}, if whenever \(A,B,C\in\mathcal K\) and
    \[
        e\colon A\hookrightarrow B,
        \qquad
        f\colon A\hookrightarrow C
    \]
    are embeddings, there exist \(D\in\mathcal K\) and embeddings
    \[
        g\colon B\hookrightarrow D,
        \qquad
        h\colon C\hookrightarrow D
    \]
    such that
    \[
        g\circ e=h\circ f.
    \]
\end{definition}

Thus \emph{(AP)} says that any two structures in \(\mathcal K\) which share a common substructure can be embedded into a larger structure in \(\mathcal K\) so that the two copies of the common part coincide.

\[
    \begin{tikzcd}[row sep=3em, column sep=3.5em, ampersand replacement=\&]
        \& B \ar[dashed,dr,"g"] \& \\
        A \ar[ur,"e"] \ar[dr,"f"']
        \& \& D \\
        \& C \ar[dashed,ur,"h"'] \&
    \end{tikzcd}
\]

We now return to the class of finite linear orders.

\begin{proposition}\label{prop:finite-linear-orders-ap}
    The class of finite linear orders has the amalgamation property.
\end{proposition}

\begin{proof}
    We first treat the special case in which \(A\subseteq B\) and \(A\subseteq C\), the maps \(e\) and \(f\) are the inclusion maps, and
    \[
        A=B\cap C.
    \]
    Write
    \[
        A=\{a_1<\cdots<a_n\}.
    \]
    For convenience, introduce formal symbols
    \[
        a_0=-\infty,\qquad a_{n+1}=+\infty.
    \]
    For each \(0\leq i\leq n\), let
    \[
        B_i:=\{b\in B\setminus A : a_i<b<a_{i+1}\},
        \qquad
        C_i:=\{c\in C\setminus A : a_i<c<a_{i+1}\},
    \]
    where the inequalities involving \(a_0\) and \(a_{n+1}\) are interpreted in the obvious way. Since \(A\) is a substructure of both \(B\) and \(C\), every element of \(B\setminus A\) belongs to exactly one \(B_i\), and every element of \(C\setminus A\) belongs to exactly one \(C_i\).

    Let
    \[
        D:=B\cup C
    \]
    as a set. We define a linear order \(<_{D}\) on \(D\) as follows:
    \begin{itemize}
        \item \(<_{D}\) agrees with the given orders on \(B\) and on \(C\);
        \item for each \(i\), every element of \(B_i\) is declared smaller than every element of \(C_i\);
        \item if \(x\in B_i\cup C_i\) and \(y\in B_j\cup C_j\) with \(i<j\), then \(x<_{D}y\);
        \item if \(x\in B_i\cup C_i\), then
              \[
                  a_i <_{D} x <_{D} a_{i+1},
              \]
              with the obvious interpretation when \(i=0\) or \(i=n\).
    \end{itemize}

    In other words, inside each interval determined by the elements of \(A\), we place first the elements coming from \(B\), in their original order, and then the elements coming from \(C\), in their original order.

    It is straightforward to verify that \(<_{D}\) is a linear order on \(D\). By construction, the inclusion maps
    \[
        j_B\colon B\hookrightarrow D,
        \qquad
        j_C\colon C\hookrightarrow D
    \]
    are embeddings of linear orders, and they agree on \(A\). Thus \(D\) is an amalgam of \(B\) and \(C\) over \(A\).
    We now consider the general case. Let
    \[
        e\colon A\hookrightarrow B,
        \qquad
        f\colon A\hookrightarrow C
    \]
    be embeddings of finite linear orders. By Lemma~\ref{lem:embedding-as-inclusion}, there exist finite linear orders \(B_1\) and \(C_1\), together with isomorphisms
    \[
        \beta\colon B_1\to B,
        \qquad
        \gamma\colon C_1\to C,
    \]
    such that \(A\subseteq B_1\), \(A\subseteq C_1\), and
    \[
        \beta\!\restriction_A=e,
        \qquad
        \gamma\!\restriction_A=f.
    \]
    Replacing \(B_1\) and \(C_1\) by isomorphic copies if necessary, we may moreover assume that
    \[
        B_1\cap C_1=A.
    \]
    By the special case already proved, there exist a finite linear order \(D_1\) and embeddings
    \[
        j_B\colon B_1\hookrightarrow D_1,
        \qquad
        j_C\colon C_1\hookrightarrow D_1
    \]
    such that
    \[
        j_B\!\restriction_A=j_C\!\restriction_A.
    \]
    Define
    \[
        g:=j_B\circ \beta^{-1}\colon B\hookrightarrow D_1,
        \qquad
        h:=j_C\circ \gamma^{-1}\colon C\hookrightarrow D_1.
    \]
    Then
    \[
        g\circ e
        =
        j_B\circ \beta^{-1}\circ e
        =
        j_B\!\restriction_A
        =
        j_C\!\restriction_A
        =
        j_C\circ \gamma^{-1}\circ f
        =
        h\circ f,
    \]
    so \(D_1\) is an amalgam of \(e\) and \(f\).

    \[
        \begin{tikzcd}[column sep=large, row sep=large, ampersand replacement=\&]
            \& B_1 \arrow[rr, "\beta", "\cong"'] \arrow[dr, "j_B"'] \& \& B \arrow[dl, dashed, "g:=j_B\circ\beta^{-1}"] \\
            A \arrow[ur, hook] \arrow[dr, hook'] \arrow[urrr, bend left=5, "e"] \arrow[drrr, bend right=5, "f"'] \& \& D_1 \& \\
            \& C_1 \arrow[rr, "\gamma"', "\cong"] \arrow[ur, "j_C"] \& \& C \arrow[ul, dashed, "h:=j_C\circ\gamma^{-1}"']
        \end{tikzcd}
    \]

    Hence the class of finite linear orders has \emph{(AP)}.

\end{proof}

This example leads to the central class of finite structures considered in Fraïssé theory.

\begin{definition}[Fraïssé class]
    Let \(L\) be a countable signature. A class \(\mathcal K\) of finitely generated \(L\)-structures is called a \emph{Fraïssé class} if it satisfies the following conditions:
    \begin{enumerate}
        \item \(\mathcal K\) is closed under isomorphism;
        \item \(\mathcal K\) has countably many isomorphism types;
        \item \(\mathcal K\) has the hereditary property \emph{(HP)};
        \item \(\mathcal K\) has the joint embedding property \emph{(JEP)};
        \item \(\mathcal K\) has the amalgamation property \emph{(AP)}.
    \end{enumerate}
\end{definition}


\section{Fraïssé limits}

We now arrive at the central notion of the chapter.

\begin{definition}[Ultrahomogeneous structure]
    Let \(D\) be an \(L\)-structure. We say that \(D\) is \emph{ultrahomogeneous} if every isomorphism between finitely generated substructures of \(D\) extends to an automorphism of \(D\).
\end{definition}

Fraïssé's theorem shows that every Fraïssé class determines, up to isomorphism, a unique countable ultrahomogeneous structure whose age is exactly that class.

\begin{theorem}[Fraïssé's theorem]\label{thm:fraisse-theorem}
    Let \(\mathcal K\) be a Fraïssé class. Then there exists a countable ultrahomogeneous \(L\)-structure \(D\) such that
    \[
        \Age(D)=\mathcal K.
    \]
    Moreover, \(D\) is unique up to isomorphism.
\end{theorem}

\begin{definition}[Fraïssé limit]
    Let \(\mathcal K\) be a Fraïssé class. The unique countable ultrahomogeneous structure \(D\) satisfying
    \[
        \Age(D)=\mathcal K
    \]
    is called the \emph{Fraïssé limit} of \(\mathcal K\).
\end{definition}

The next two sections are devoted to the proof of Theorem~\ref{thm:fraisse-theorem}: first the uniqueness of the Fraïssé limit, and then its existence.

\section{Uniqueness of the Fraïssé limit}

We begin with a weaker extension property which is sufficient, in the countable case, to determine a structure uniquely from its age.

\begin{definition}[Weak homogeneity]
    Let \(D\) be an \(L\)-structure. We say that \(D\) is \emph{weakly homogeneous} if whenever \(A\subseteq B\) are finitely generated substructures of \(D\) and
    \[
        f\colon A\hookrightarrow D
    \]
    is an embedding, there exists an embedding
    \[
        g\colon B\hookrightarrow D
    \]
    extending \(f\).
\end{definition}

\[
    \begin{tikzcd}[ampersand replacement=\&]
        A \arrow[r,"f"] \arrow[d,hook] \& D \\
        B \arrow[ur,dashed,"g"']
    \end{tikzcd}
\]

Weak homogeneity is strictly weaker than ultrahomogeneity as a formal condition, but for finite or countable structures the two notions coincide.

\begin{example}
    The ordered set \((\mathbb Z,<)\) is not weakly homogeneous.
\end{example}

\begin{proof}
    Let
    \[
        A=\{0,2\}\subseteq \mathbb Z,
        \qquad
        B=\{0,1,2\}\subseteq \mathbb Z,
    \]
    with the induced order, and define \(f\colon A\to \mathbb Z\) by
    \[
        f(0)=0,
        \qquad
        f(2)=1.
    \]
    Then \(f\) is an embedding. If \(g\colon B\to \mathbb Z\) were an embedding extending \(f\), we would have
    \[
        0=g(0)<g(1)<g(2)=1,
    \]
    which is impossible in \(\mathbb Z\). Hence no such extension exists.
\end{proof}

\begin{lemma}\label{lem:ultra-implies-weak}
    Every ultrahomogeneous structure is weakly homogeneous.
\end{lemma}

\begin{proof}
    Let \(A\subseteq B\) be finitely generated substructures of an ultrahomogeneous structure \(D\), and let \(f\colon A\hookrightarrow D\) be an embedding. Then \(f\colon A\to f(A)\) is an isomorphism between finitely generated substructures of \(D\). By ultrahomogeneity, \(f\) extends to an automorphism \(\alpha\) of \(D\). The restriction
    \[
        g:=\alpha\!\restriction_B\colon B\hookrightarrow D
    \]
    is then an embedding extending \(f\).
\end{proof}

The next lemma shows that a countable weakly homogeneous structure is universal for its age.

\begin{lemma}\label{lem:weak-age-embedding}
    Let \(C\) and \(D\) be at most countable \(L\)-structures. Suppose that
    \[
        \mathrm{Age}(C)\subseteq \mathrm{Age}(D)
    \]
    and that \(D\) is weakly homogeneous. Then every embedding of a finitely generated substructure of \(C\) into \(D\) extends to an embedding of \(C\) into \(D\).
\end{lemma}

\begin{proof}
    Let \(f_0\colon A_0\hookrightarrow D\) be an embedding, where \(A_0\subseteq C\) is finitely generated. Choose an increasing chain of finitely generated substructures
    \[
        A_0\subseteq A_1\subseteq A_2\subseteq \cdots
    \]
    with
    \[
        C=\bigcup_{n<\omega} A_n .
    \]
    For example, if \(C\setminus A_0=\{c_0,c_1,\dots\}\), we may take
    \[
        A_{n+1}:=\langle A_n\cup\{c_n\}\rangle .
    \]

    We construct inductively embeddings
    \[
        f_n\colon A_n\hookrightarrow D
        \qquad (n<\omega)
    \]
    such that \(f_n\subseteq f_{n+1}\). Suppose \(f_n\) has been constructed. Since \(A_{n+1}\) is finitely generated and \(\mathrm{Age}(C)\subseteq \mathrm{Age}(D)\), there exist a finitely generated substructure \(B\subseteq D\) and an isomorphism
    \[
        u\colon A_{n+1}\xrightarrow{\ \cong\ } B.
    \]
    Then
    \[
        f_n\circ (u\!\restriction_{A_n})^{-1}\colon u(A_n)\hookrightarrow D
    \]
    is an embedding. Since \(u(A_n)\subseteq B\) and \(D\) is weakly homogeneous, this embedding extends to an embedding
    \[
        h\colon B\hookrightarrow D.
    \]
    Define
    \[
        f_{n+1}:=h\circ u\colon A_{n+1}\hookrightarrow D.
    \]
    For \(a\in A_n\),
    \[
        f_{n+1}(a)=h(u(a))
        =\bigl(f_n\circ (u\!\restriction_{A_n})^{-1}\bigr)(u(a))
        =f_n(a),
    \]
    so \(f_{n+1}\) extends \(f_n\).

    Now let
    \[
        f:=\bigcup_{n<\omega} f_n\colon C\to D.
    \]
    This is well defined because the maps \(f_n\) are compatible. Since every finite tuple from \(C\) is contained in some \(A_n\), and each \(f_n\) is an embedding, it follows that \(f\) is an embedding of \(C\) into \(D\). By construction, \(f\) extends \(f_0\).
\end{proof}

We now obtain the uniqueness statement by the usual back-and-forth argument.

\begin{lemma}\label{lem:countable-weak-uniqueness}
    Let \(C\) and \(D\) be at most countable weakly homogeneous \(L\)-structures with
    \[
        \mathrm{Age}(C)=\mathrm{Age}(D).
    \]
    Then \(C\cong D\). More precisely, every isomorphism between finitely generated substructures of \(C\) and \(D\) extends to an isomorphism from \(C\) onto \(D\).
\end{lemma}

\begin{proof}
    Let
    \[
        f_0\colon A_0\xrightarrow{\ \cong\ } B_0
    \]
    be an isomorphism between finitely generated substructures \(A_0\subseteq C\) and \(B_0\subseteq D\). Choose increasing chains of finitely generated substructures
    \[
        A_0\subseteq A_1\subseteq A_2\subseteq \cdots,
        \qquad
        B_0\subseteq B_1\subseteq B_2\subseteq \cdots
    \]
    such that
    \[
        C=\bigcup_{n<\omega} A_n,
        \qquad
        D=\bigcup_{n<\omega} B_n .
    \]

    We construct inductively a chain of partial isomorphisms
    \[
        f_0\subseteq f_1\subseteq f_2\subseteq \cdots
    \]
    between finitely generated substructures of \(C\) and \(D\) such that for every \(n\),
    \[
        A_n\subseteq \dom(f_{2n}),
        \qquad
        B_n\subseteq \operatorname{im}(f_{2n+1}).
    \]

    Suppose \(f_{2n}\) has been constructed. Since \(D\) is weakly homogeneous and \(\mathrm{Age}(C)\subseteq \mathrm{Age}(D)\), Lemma~\ref{lem:weak-age-embedding} extends \(f_{2n}\) to an embedding
    \[
        g\colon C\hookrightarrow D.
    \]
    Let
    \[
        f_{2n+1}:=g\!\restriction_{\langle \dom(f_{2n})\cup A_{n+1}\rangle}.
    \]
    Then \(f_{2n+1}\) extends \(f_{2n}\), and its domain contains \(A_{n+1}\).

    Now \(f_{2n+1}^{-1}\) is an isomorphism from the finitely generated substructure \(\operatorname{im}(f_{2n+1})\subseteq D\) into \(C\). Since \(C\) is weakly homogeneous and \(\mathrm{Age}(D)\subseteq \mathrm{Age}(C)\), Lemma~\ref{lem:weak-age-embedding} extends \(f_{2n+1}^{-1}\) to an embedding
    \[
        h\colon D\hookrightarrow C.
    \]
    Define
    \[
        f_{2n+2}:=\bigl(h\!\restriction_{\langle \operatorname{im}(f_{2n+1})\cup B_{n+1}\rangle}\bigr)^{-1}.
    \]
    Then \(f_{2n+2}\) extends \(f_{2n+1}\), and its image contains \(B_{n+1}\).

    Thus the construction continues indefinitely. Let
    \[
        f:=\bigcup_{n<\omega} f_n.
    \]
    Because the maps \(f_n\) are compatible, \(f\) is a well-defined embedding \(C\to D\). The domain condition implies that every element of \(C\) lies in \(\dom(f)\), and the image condition implies that every element of \(D\) lies in \(\operatorname{im}(f)\). Hence \(f\) is an isomorphism from \(C\) onto \(D\). By construction, it extends \(f_0\).
\end{proof}

\begin{corollary}\label{cor:weak-equals-ultra-countable}
    For a finite or countable \(L\)-structure \(D\), the following are equivalent:
    \begin{enumerate}
        \item \(D\) is ultrahomogeneous;
        \item \(D\) is weakly homogeneous.
    \end{enumerate}
\end{corollary}

\begin{proof}
    By Lemma~\ref{lem:ultra-implies-weak}, every ultrahomogeneous structure is weakly homogeneous.

    Conversely, suppose that \(D\) is finite or countable and weakly homogeneous. Let
    \[
        f\colon A\xrightarrow{\ \cong\ } B
    \]
    be an isomorphism between finitely generated substructures of \(D\). Applying Lemma~\ref{lem:countable-weak-uniqueness} with \(C=D\), \(D=D\), and the given partial isomorphism \(f\), we obtain an automorphism of \(D\) extending \(f\). Hence \(D\) is ultrahomogeneous.
\end{proof}

In particular, if \(C\) and \(D\) are countable ultrahomogeneous structures with the same age, then \(C\cong D\). This proves the uniqueness part of Fraïssé's theorem.



\section{Existence of the Fraïssé limit}

We now prove the existence part of Fraïssé's theorem. The idea is to construct a countable structure \(D\) with age \(\mathcal K\) and with the extension property that characterizes weak homogeneity.

We first record a simple lemma on increasing unions.

\begin{lemma}\label{lem:age-of-union}
    Let
    \[
        D_0\subseteq D_1\subseteq D_2\subseteq \cdots
    \]
    be an increasing chain of \(L\)-structures, and let
    \[
        D:=\bigcup_{i<\omega} D_i.
    \]
    Let \(\mathcal K\) be a class of finitely generated \(L\)-structures.
    \begin{enumerate}
        \item[(i)] If \(\mathrm{Age}(D_i)\subseteq \mathcal K\) for all \(i<\omega\), then
              \[
                  \mathrm{Age}(D)\subseteq \mathcal K.
              \]
        \item[(ii)] If, in addition, every \(A\in\mathcal K\) embeds into some \(D_i\), then
              \[
                  \mathrm{Age}(D)=\mathcal K.
              \]
    \end{enumerate}
\end{lemma}

\begin{proof}
    Let \(A\) be a finitely generated substructure of \(D\), generated by a finite tuple \(\bar a\). Since the chain is increasing, all entries of \(\bar a\) lie in some \(D_i\). Hence
    \[
        A=\langle \bar a\rangle_D=\langle \bar a\rangle_{D_i},
    \]
    so \(A\in\mathrm{Age}(D_i)\subseteq\mathcal K\). This proves (i).

    For (ii), part (i) already gives \(\mathrm{Age}(D)\subseteq\mathcal K\). Conversely, let \(A\in\mathcal K\). By assumption, \(A\) embeds into some \(D_i\), and therefore also into \(D\). Thus \(A\in\mathrm{Age}(D)\), so \(\mathcal K\subseteq\mathrm{Age}(D)\).
\end{proof}

\begin{proof}[Proof of Fraïssé's theorem]
    Let \(\mathcal K\) be a Fraïssé class. Since only isomorphism types matter, we may assume that \(\mathcal K\) is closed under isomorphism.

    We shall construct an increasing chain
    \[
        D_0\subseteq D_1\subseteq D_2\subseteq \cdots
    \]
    of members of \(\mathcal K\) such that the following extension property holds:
    \begin{quote}
        whenever \(A,B\in\mathcal K\) with \(A\subseteq B\), and \(f\colon A\hookrightarrow D_i\) is an embedding, there exists \(j>i\) and an embedding
        \[
            g\colon B\hookrightarrow D_j
        \]
        extending \(f\).
    \end{quote}

    Assume for the moment that such a chain has been constructed, and let
    \[
        D:=\bigcup_{i<\omega} D_i.
    \]

    Since \(\mathcal K\) has countably many isomorphism types and each member of \(\mathcal K\) is finitely generated in a countable signature, each \(D_i\) is countable. Hence \(D\) is countable.

    Because each \(D_i\) belongs to \(\mathcal K\), Lemma~\ref{lem:age-of-union}(i) gives
    \[
        \mathrm{Age}(D)\subseteq \mathcal K.
    \]

    We claim that in fact \(\mathrm{Age}(D)=\mathcal K\). Let \(A\in\mathcal K\). By \emph{(JEP)}, there exist \(B\in\mathcal K\) and embeddings
    \[
        u\colon D_0\hookrightarrow B,
        \qquad
        v\colon A\hookrightarrow B.
    \]
    By Lemma~\ref{lem:embedding-as-inclusion}, after replacing \(B\) by an isomorphic copy, we may assume that
    \[
        D_0\subseteq B.
    \]
    Now apply the extension property to the identity embedding \(1_{D_0}\colon D_0\hookrightarrow D_0\). Then for some \(j>0\), the inclusion \(D_0\subseteq B\) extends to an embedding
    \[
        g\colon B\hookrightarrow D_j.
    \]
    Composing with \(v\), we obtain an embedding of \(A\) into \(D_j\). Thus every \(A\in\mathcal K\) embeds into some \(D_j\), and Lemma~\ref{lem:age-of-union}(ii) yields
    \[
        \mathrm{Age}(D)=\mathcal K.
    \]

    We next show that \(D\) is weakly homogeneous. Let \(A\subseteq B\) be finitely generated substructures of \(D\), and let
    \[
        f\colon A\hookrightarrow D
    \]
    be an embedding. Since \(A\) is finitely generated and \(f(A)\) is finite or countable, there exists \(i<\omega\) such that
    \[
        f(A)\subseteq D_i.
    \]
    Because \(\mathrm{Age}(D)=\mathcal K\), both \(A\) and \(B\) belong to \(\mathcal K\). Applying the extension property to the embedding
    \[
        f\colon A\hookrightarrow D_i,
    \]
    we obtain \(j>i\) and an embedding
    \[
        g\colon B\hookrightarrow D_j\subseteq D
    \]
    extending \(f\). Hence \(D\) is weakly homogeneous.

    Since \(D\) is countable, Corollary~\ref{cor:weak-equals-ultra-countable} implies that \(D\) is ultrahomogeneous. Therefore \(D\) is a countable ultrahomogeneous structure with age \(\mathcal K\), as required.

    It remains to construct the chain \((D_i)_{i<\omega}\).

    Let
    \[
        \mathcal P=\{(A_n,B_n): n<\omega\}
    \]
    be a countable set of representatives of the isomorphism types of pairs \((A,B)\) such that \(A,B\in\mathcal K\) and \(A\subseteq B\). Fix a bijection
    \[
        \pi\colon \omega\times\omega\to\omega
    \]
    such that \(\pi(i,j)\geq i\) for all \(i,j<\omega\). Choose any \(D_0\in\mathcal K\).

    Suppose \(D_i\) has been constructed. Since \(D_i\) is countable and each \(A_n\) is finitely generated, there are only countably many embeddings of the various \(A_n\) into \(D_i\). Enumerate all triples
    \[
        (f_{ij},A_{ij},B_{ij}) \qquad (j<\omega),
    \]
    where \(A_{ij}\subseteq B_{ij}\) is one of the pairs from \(\mathcal P\) and
    \[
        f_{ij}\colon A_{ij}\hookrightarrow D_i
    \]
    is an embedding.

    Now suppose \(D_k\) has been defined, and write
    \[
        k=\pi(i,j).
    \]
    Then
    \[
        f_{ij}\colon A_{ij}\hookrightarrow D_i\subseteq D_k
    \]
    is an embedding. Apply \emph{(AP)} to the inclusion
    \[
        A_{ij}\hookrightarrow B_{ij}
    \]
    and the embedding
    \[
        f_{ij}\colon A_{ij}\hookrightarrow D_k.
    \]
    We obtain some \(E\in\mathcal K\) and embeddings
    \[
        r\colon B_{ij}\hookrightarrow E,
        \qquad
        s\colon D_k\hookrightarrow E
    \]
    such that
    \[
        r\!\restriction_{A_{ij}}=s\circ f_{ij}.
    \]
    By Lemma~\ref{lem:embedding-as-inclusion}, after replacing \(E\) by an isomorphic copy, we may assume that
    \[
        D_k\subseteq E
    \]
    and that \(s\) is the inclusion map. Set
    \[
        D_{k+1}:=E.
    \]
    Then the embedding \(r\colon B_{ij}\hookrightarrow D_{k+1}\) extends \(f_{ij}\).

    This completes the construction of the chain. Because every embedding problem over every stage \(D_i\) appears as some triple \((f_{ij},A_{ij},B_{ij})\), and because \(\pi(i,j)\geq i\), each such problem is eventually handled at a later stage. Hence the required extension property holds.
\end{proof}


\section{The converse direction}

We now show that the conditions in Fraïssé's theorem are also necessary.

\begin{theorem}\label{thm:converse-fraisse}
    Let \(L\) be a countable signature, and let \(D\) be a finite or countable ultrahomogeneous \(L\)-structure. Then \(\mathrm{Age}(D)\) is a Fraïssé class. More explicitly, \(\mathrm{Age}(D)\) is nonempty, has countably many isomorphism types, and satisfies \emph{(HP)}, \emph{(JEP)}, and \emph{(AP)}.
\end{theorem}

\begin{proof}
    Set
    \[
        \mathcal K:=\mathrm{Age}(D).
    \]
    By Proposition~\ref{prop:age-hp-jep}, the class \(\mathcal K\) is nonempty and satisfies \emph{(HP)} and \emph{(JEP)}.

    It remains to verify that \(\mathcal K\) has countably many isomorphism types and satisfies \emph{(AP)}.

    Since \(D\) is finite or countable, there are only countably many finite tuples from \(D\). Every member of \(\mathcal K\) is finitely generated, hence is generated by some finite tuple from \(D\). Therefore \(\mathcal K\) has at most countably many isomorphism types.

    To prove \emph{(AP)}, let \(A,B,C\in\mathcal K\), and let
    \[
        e\colon A\hookrightarrow B,
        \qquad
        f\colon A\hookrightarrow C
    \]
    be embeddings. Since \(A\in\mathcal K=\mathrm{Age}(D)\), there exists an embedding
    \[
        u\colon A\hookrightarrow D.
    \]
    By Corollary~\ref{cor:weak-equals-ultra-countable}, the structure \(D\) is weakly homogeneous. Now
    \[
        u\circ e^{-1}\colon e(A)\hookrightarrow D
    \]
    is an embedding of the finitely generated substructure \(e(A)\subseteq B\) into \(D\). Since \(B\in\mathcal K\), weak homogeneity yields an embedding
    \[
        j_B\colon B\hookrightarrow D
    \]
    extending \(u\circ e^{-1}\). Thus
    \[
        j_B\circ e=u.
    \]
    Similarly, the embedding
    \[
        u\circ f^{-1}\colon f(A)\hookrightarrow D
    \]
    extends to an embedding
    \[
        j_C\colon C\hookrightarrow D
    \]
    with
    \[
        j_C\circ f=u.
    \]
    Hence
    \[
        j_B\circ e=j_C\circ f.
    \]

    Let
    \[
        E:=\langle j_B(B)\cup j_C(C)\rangle \subseteq D.
    \]
    Since \(B\) and \(C\) are finitely generated, the structure \(E\) is finitely generated, so \(E\in\mathcal K\). The maps \(j_B\colon B\hookrightarrow E\) and \(j_C\colon C\hookrightarrow E\) are embeddings and satisfy
    \[
        j_B\circ e=j_C\circ f,
    \]
    so \(E\) is an amalgam of \(B\) and \(C\) over \(A\). Therefore \(\mathcal K\) satisfies \emph{(AP)}.

    Thus \(\mathcal K\) is a Fraïssé class.
\end{proof}